\begin{annexitem}{annex:p12.n-c-035.5.minimal-sets}{thm:p12-eu-minimal-sets}{Comparison and minimal classes}
\begin{proof}
Reflexivity uses identity. If \(x\theta=_E x\sigma\lambda\) and
\(x\rho=_E x\theta\mu\), substitution closure gives
\(x\rho=_E x\sigma(\lambda\mu)\). Term composition follows by
induction on finite terms; composition still has finite support.
Thus the relation is transitive. Its symmetric part is an equivalence,
and the quotient comparison is well defined and antisymmetric.

In any partial order, a complete subset is inclusion-minimal exactly
when it is an antichain. If it contains \(c<d\), deleting \(d\)
retains coverage, since anything above \(d\) is above \(c\).
Conversely, if a proper subset still covers, a deleted \(d\) is above
some retained \(c\), exhibiting two distinct comparable members.
For a complete antichain \(C\), every minimal element \(m\) is
covered by \(c\leq m\), forcing \(c=m\). If \(d\in C\) were
not minimal, choose \(u<d\) and \(c\in C\) with \(c\leq u\),
contradicting the antichain property. Therefore \(C\) is exactly
the set of minimal elements. Conversely, if that set is complete,
it is an antichain and hence minimal complete. On lifting back from
classes, completeness requires at least one representative of each
minimal class, and incomparability permits at most one. This proves
existence, cardinality, and the stated uniqueness; it also handles
the empty solution set without creating a spurious unifier.
\end{proof}
\end{annexitem}

\begin{annexitem}{annex:p12.n-c-035.5.ordered-closure}{thm:p12-eu-ordered-closure}{The ordered overlap tests, direct convergence, and the comparison invariant}
\begin{proof}
Every admitted pair is an E-consequence, by substituting its syntactic
overlap unifier in the two component identities. Induction on levels
therefore preserves exactly the original congruence. Every oriented
ground instance decreases the fixed well-founded order, so termination
holds even for the infinite closure.
On ground terms this is exactly the source's context-ordered rewrite
relation: totality compares the two local replacement terms, and
monotonicity makes that direction agree with their enclosing context.
Hence using decreasing ground instances has not silently enlarged
the source rewrite relation.

First handle possible variable-left ground rules. If the closure
contains \(x\approx r\) with x absent from r, and r is another
variable, it already contains a two-distinct-variable identity. If r
is nonvariable, overlap the reverse identity \(r\approx x\)
with its fresh variant \(r'\approx y\) at the root. A syntactic
MGU of the two renamed r-patterns exists. It fixes the independent
right variables x,y. Its common term \(r\mu\) contains neither.
A term v not containing x is incomparable with x under any strict
substitution-stable order: either \(v\succ x\) or \(x\succ v\)
would become \(v\succ v\) on assigning x to v. Thus the identity
witness satisfies both source not-\(\preceq\) tests, and the
ordered critical identity \(x\approx y\) is added at the next
level. Its ground instances compare every pair of ground terms;
all reduce to the least ground term of the well-founded total order.
Hence ground confluence follows in this case, including trivial theories.

In the remaining case, every variable-left identity has that variable
in its other side. Apart from an identical variable equation, any
ground instance of the other side properly contains the variable's
ground image. Monotonicity and well-foundedness forbid reducing a
term to a proper context containing it: iterating that context would
give an infinite descent. Thus all actual decreasing ground rules
have nonvariable pattern left sides.

Induct on a ground source w to prove unique normal forms. All its
proper reducts are smaller and already normalize uniquely. Two first
steps at disjoint occurrences commute. For a nonvariable nested
overlap, take the fresh component identities from their finite
closure levels and a syntactic overlap MGU. The actual ground
matching substitution factors through that MGU. It witnesses both
ordered tests at the local common source: the outer ground rule is
oriented downward, and the inner ground rule decreases through its
inner context by monotonicity. Thus the local source is strictly
greater than each local branch, before adding the common outer context.
Consequently this critical identity is admitted at a later finite
level. Its two instantiated branches, in their context, are either
identical or related by one oriented ground step, so their already
unique smaller normal forms agree.

For an overlap inside a variable image of the outer pattern, update
that image by the inner rewrite in every other left-side occurrence
and in every right-side occurrence. Both peak branches are thereby
connected to the two sides of the updated outer identity. Every
intermediate term is strictly below w; that updated identity is
already in the closure and can be used in whichever ground direction
decreases the order. This covers duplicated and erased occurrences
and does not assume its former orientation survives the image update.
The smaller-term induction makes every neighboring rewrite pair in
this conversion share a normal form. Thus the two branches agree.
All first steps from w agree, proving its unique normal form and
ground confluence. This is a direct admission/peak proof for the
ordered closure, not invocation of unrestricted closure completeness.

Ground equational conversions of E can be grounded at intermediate
variables, fixing their endpoints, and then oriented according to
the total order. Thus E-equal ground terms have the same ordered
normal form, and conversely every ordered reduction is an E-step.
For nontrivial E the earlier unrestricted theorem applies as well.
Let m be the least ground member of a given E-class. If either
normal form n were greater than m, confluence and equivalence would
require m to reduce to n, contradicting strict decrease. Therefore
both normal forms equal m. This establishes the genuine comparison
without identifying equation sets or rewrite steps.

For transporting the two-direction critical simulation to the
source G proof, choose a complete simplification ordering with the
proper-subterm property also on open terms, for example LPO.
In a nontrivial theory a variable-left inner identity \(x\approx r\)
must contain x in r. If r is x the replacement branch equals the
source; otherwise the replaced subterm is a proper subterm of its
replacement, and monotonicity lifts that strict increase through
the unchanged surrounding context, for every substitution. The second ordered
test therefore excludes this overlap. All admitted critical pairs
then have the two nonvariable left patterns of the earlier complete
simulation. With the same fresh-variable and actual-U selector,
\(E_j\subseteq F_j\) follows inductively, but only under these
declared restrictions; it is not asserted for every weaker open
order or every selector convention.

One may now run the already constructed innermost ground-witness G
lift using the directly proved ordered \(E_\omega\) normalization:
the reduced-image/nonvariable-position proof and the mandatory
matching-decomposition multiset descent use only its actual decreasing
ground steps. Every used identity has a finite ordered-closure depth,
and the proved two-direction simulation eliminates its parent
critical identities with strictly decreasing depth multiset. Hence
this is a genuine ordered-source route back to E, alongside the
separately labelled unrestricted route. No test decidability or
finite completion follows from these existence constructions.
\end{proof}
\end{annexitem}

\begin{annexitem}{annex:p12.n-c-035.5.local-simplifications}{thm:p12-eu-local-simplifications}{Local witness bijections and each actual descent coordinate}
\begin{proof}
Use a ground reduced witness after the declared parameter freeze.
For Problem Decomposition, the root f cannot be rewritten by R.
Normalizing its arguments therefore yields a normal form still
rooted at f. Two such terms are E-equal exactly when corresponding
argument normal forms agree, by confluence and literal roots.
Thus the parent E-equation and its child E-equations have exactly
the same solutions. Replacing each instantiated parent by its
proper arguments strictly decreases M, including arity zero.

For corrected Propagation, literal solution of S gives
\(\sigma_S\rho=\rho\). The side condition consequently gives
\(u\rho=s\rho\). Every old solution extends uniquely at fresh y
by \(y\rho=s\rho\), and the replacement equation has literally
the same instantiated terms. Conversely every new solution restricts
to an old one by its retained y-binding. If the old witness is
reduced, the solved left variable's image \(x\rho=t[s]\rho\)
is normal; its subterm \(s\rho\) is normal. The fresh y-image
is therefore reduced. M is unchanged. Replacing a nonvariable u
by a variable removes at least one nonvariable skeleton symbol,
so i strictly decreases. If u is already a variable this last
claim is false; that is precisely why the printed literal variant
is not used in the checked simplification class.

For cleaned Reduce, its fresh-only matching map constrains no old
variable and satisfies \(l\gamma=u\sigma_S\). Under every old
literal S-solution this becomes \(l\gamma\rho=u\rho\).
Replacing u by \(r\gamma\) is therefore an actual R step on
the instantiated equation term and an E-conversion for every old
solution. Both implications for equation solutions hold, while
C,S are unchanged and all old reduced images stay unchanged.
No fresh rule variable remains in \(r\gamma\): its variables
come from the matched old term, since R has no right-only variable.
All discarded fresh match domains are absent from old data and
input variables. Dropping their isolated solved bindings therefore
loses no old solution; they can always be restored as their forced
literal values if the literal source state is desired. The changed
ground skeleton term strictly decreases under R, hence M decreases.

The syntactic C rules retain the literal witness and solved-map
calculation already proved. For B they leave skeleton size unchanged:
elimination decreases n1; deletion/decomposition decrease n2 if
n1 stays fixed; orientation decreases n3 if both stay fixed.
For N elimination may grow i, but its n1 decrease precedes i in
the declared N tuple. Other C rules cannot grow that last coordinate.
Propagation leaves C unchanged and lowers i in either tuple; the
two equation simplifications lower M before every other coordinate.
Thus every actual member of the checked class has its asserted
strict well-founded descent and reduced-witness preservation.
\end{proof}
\end{annexitem}

\begin{annexitem}{annex:p12.n-c-035.5.eager-corrected}{thm:p12-eu-eager-corrected}{Eager simplification as an induction, not a completeness assumption}
\begin{proof}
Take the normalized ground witness of an arbitrary original unifier.
At every state where a checked simplification applies, choose one.
The local proof extends or retains the reduced witness without
changing its old-variable assignment, and strictly decreases the
specified well-founded tuple. Thus this eager phase is finite;
it cannot prevent a completion step by an infinite fresh-alias loop.
When none applies, use one of the actual complete B/N cases:
solve a literal C constraint, Constrain an equation whose ground
sides are the same normal term, or lift an innermost ground R step.
The syntactic rules are already in the eager class if C is nonempty.
Constrain removes its two skeleton terms and hence decreases M.
The lifted LP block replaces a ground term by its strict reduct;
fresh rule images are normal proper subterms of the chosen innermost
redex, and its mandatory first matching Decomposition is performed.
Thus it decreases M and preserves the reduced witness, just as in
the independently proved B/N lifting cases.

Repeat induction on the tuple. At the terminal empty P,C, the
literal S-solution gives factorization through \(\sigma_S\).
Every fresh witness extension was explicitly fixed at a normal
old subterm and every cleaned match projection involved only fresh
domains absent from original data. Restrict to original variables
and unfreeze the finite parameter constants to obtain coverage of
the arbitrary original E-unifier. Thus eager control is proved
from checked local operations and the real lifting cases; it is
not supplied by an assumption that the desired strategy is complete.
Backward soundness also holds for each checked operation, so every
successful output is an original E-unifier.

For the source abstract class in Definition 4.15, the same induction
works for any actual rule for which its reduced-witness equivalence
and strict-descent conditions have been proved. Those local conditions
must still be checked; the literal aliasing Propagation and uncleaned
fresh-binding certificate are not silently included in that class.
\end{proof}
\end{annexitem}

\begin{annexitem}{annex:p12.n-c-035.5.eager-reduce-literal}{thm:p12-eu-eager-reduce-literal}{Eager literal Reduce with explicitly inactive match domains}
\begin{proof}
Separate the fresh match domains retained by completed Reduce steps
from active variables. Such a domain occurs only as a left side of
its own solved equation: it is absent from P,C, all other solved
right sides and the original variable set. Its own image uses only
older variables, because matching assigned fresh rule variables to
the old \(u\sigma_S\), and the fully substituted \(r\gamma\)
contains no fresh match variable. Later constraint substitutions
update this isolated right side but do not introduce its left
variable into any active datum. New variants are fresh for it too.
This separation is permanent.

Maintain a reduced witness on all active variables, with each
inactive match-domain value set to its forced literal solved image.
Those latter values need not be reduced. Each Reduce leaves the
old witness unchanged, adds only such inactive values and changes
one ground skeleton term by a strict R reduction, so M decreases.
An eager Reduce phase must therefore be finite. When none applies,
the B/N induction cases still apply to active data: each variable
occurring in P has a reduced image, so a needed ground redex lifts
at a nonvariable skeleton position; each C equation has a literal
witness and contains no inactive match domain. Its syntactic solving
does not constrain those isolated left variables. Constrain and LP
have their former descent and fresh-extension calculations.

The multiset/constraint induction reaches empty P,C. The full
terminal S is literally solved, including the forced inactive
images. Its factorization covers the original variables, which
are active and were never placed among the irrelevant fresh
domains. Ungrounding produces the claimed original-variable E
comparison. This proves literal eager Reduce alone. Treating all
of these fresh solved images as normal, propagating them as normal
certificates, or rejecting them by a global reduced-witness Blocking
test would invalidate this invariant and is explicitly excluded.

In the body's certificate counterexample, f and g rules terminate
by the positive interpretation f(v)=v+1, g(v)=v+1 and constants
of value one. Their only nontrivial nested interaction is inside
f's variable image and joins at a, so they are confluent.
The forced fresh z-image g(b) is nevertheless reducible. In the
literal combined branch, abstraction of the remaining g(y) into
a bound variable hides its basic redex; literal Constrain leads
to g(y)=c, a root clash. Narrowing first changes y to b and then
makes the retained z-image a syntactic redex for Blocking. These
checks establish the stated failure of the stronger certificate
combination, not failure of the sound literal Reduce operation.
\end{proof}
\end{annexitem}

\begin{annexitem}{annex:p12.n-c-035.5.ordered-narrow-blocking}{thm:p12-eu-ordered-narrow-blocking}{Every protected-position and lowest-index lifting case}
\begin{proof}
Keep the reduced ground witness and the corrected simplification
invariant. If an instantiated equation is not solved, choose a
\(\prec_{\rm red}\)-minimal redex. It is innermost: any redex
strictly inside it would be earlier by the proper-subterm priority.
No redex occurs at any earlier position, and none in any descendant
of that earlier position, since such a descendant is earlier still.
Therefore every retained earlier context region is a normal ground
term. A position inside a substituted variable image cannot be the
chosen redex, because that image is reduced.

Choose the lowest-index rule matching the ground redex and lift
the actual LP/first-Decomposition/constraint-solver block. The new
rule-variable images are normal proper subterms of the innermost
redex. Its syntactic MGU \(\nu\) factors the ground witness;
the old-variable assignments are retained by the factor. After
this block, every prescribed earlier retained context region is
still normal under that factor: it is disjoint from the replaced
hole, and its constraint-solving instances are the same ground
region. Assign a fresh abstraction variable exactly that ground
term. This is a reduced extension and literally solves its new
constraint. Abstract each original retained position once, or
abstract its maximal earlier region so that all removed descendant
positions are covered at once. Removed parts inside the old redex
are not spuriously treated as surviving positions of the context.
There are finitely many original positions. Repeated abstraction
of a newly created variable is neither required nor permitted by
the stated finite pass. Constraint solving is finite by the actual
U helper. The overall Narrow block already strictly decreased M;
the finite abstraction part changes no instantiated term.

For the lower-index test, let p=\(l\sigma_S\nu\). If p were
a literal instance of a lower-index left side, applying the ground
factor of \(\nu\) would make that lower rule match the very
same chosen ground redex. This contradicts its lowest-index choice.
Thus the actual source matching test passes on the constructed
branch. It is not necessary to replace it by E-unifiability or
to assume all branches pass.

Maintain also that every proper subterm of a C-term is normal
under the witness. Initially C is empty. A Constrain step is
chosen only for two equal normal terms; its new C-terms are
normal. An LP match consists of the innermost redex on both
sides; its proper subterms are normal even before the mandatory
first Decomposition. After that decomposition the new argument
constraints themselves are normal. Syntactic C transformations
and literal solved-image equality preserve this assertion.
The abstraction constraints have normal whole-term witnesses,
so preserve it too. Problem Decomposition, cleaned Reduce and
corrected Propagation affect no existing C-terms.

Every S-binding's instantiated right side equals its reduced
left-variable image. Hence a syntactic R-redex in an S-term
would persist under the witness and contradict normality.
Similarly a syntactic redex properly inside a C-term would
persist at that proper position, contradicting the protected
proper-subterm invariant. Eager Blocking consequently never
rejects this constructed branch, including the intermediate
matching stage of the composite rule. These are actual normality
arguments, not an assumption that Blocking is admissible.

The corrected eager-simplification induction with these finite
redex-ordered blocks therefore terminates its selected successful
branch and returns a factor of the arbitrary original witness.
Restriction to original variables and ungrounding complete the
coverage proof. Arbitrary other branches may fail or loop; neither
total termination of the search tree nor a statement that every
blocked state is E-unsatisfiable is asserted.
\end{proof}
\end{annexitem}

\begin{annexitem}{annex:p12.n-c-035.5.u-bridge}{thm:p12-eu-u-bridge}{The source's syntactic helper and its exact solved-binding state}
\begin{proof}
Use the existing composition and one-binding calculations in
\cref{thm:p12-ue-forcing}. For a source U state P;S, treat all
its equations as literal constraints. Delete and Decompose preserve
solutions by free-term identity and injective roots; Orient uses
symmetry; Clash and proper Occurs have none. At removing Elimination,
a solution satisfies \(\delta\rho=\rho\), so updates of both
P and old S preserve precisely the solutions once the binding is
retained in S. Conversely the retained binding forces the same
calculation, restoring the predecessor. Thus each actual source
transition preserves the combined literal solution set.

Inductively let D be the variables already eliminated. They occur
only as left sides in S, its right sides and P contain no member
of D, and no new variable is introduced. At a new elimination
x is in P and hence outside D; its right side has no x. Updating
old S-right-sides and P by \(\delta\) removes x everywhere
except its new single left binding. This proves the invariant
and idempotence of \(\sigma_S\). It is precisely the removing-
binding accumulator invariant: old images receive \(\delta\)
first on their terms, and the new left variable receives t, so
the accumulator update is \(\sigma_S\delta\), not its reverse.
It differs from the retained-equation sheet, which has no such S.

For termination use the source triple: distinct variables in P,
total term nodes in P counted with equation multiplicity, and bad
orientations in P. Elimination removes x from all P and introduces
no new variable, so the first coordinate decreases. Other successful
rules do not increase it. Trivial and Decomposition strictly
decrease the second coordinate if the first stays fixed; zero-arity
Decomposition removes two nodes. Orient preserves the first two
and decreases the last. Term growth after elimination is irrelevant
because its first coordinate decreases. An equation in nonempty P
has one of these six syntactic shapes, so a nonfailure terminal
state has P empty. Its \(\sigma_S\) solves S and hence the input.
Any input solution \(\rho\) literally solves S; on its left
variables and all other variables respectively the solved-form
calculation gives \(\sigma_S\rho=\rho\). This proves global
factorization, including arbitrary unused-variable images. Failure
is equivalent to absence of a solution by the invariant and
termination. Thus every chosen successful U run provides the
actual finite idempotent derivation needed by the critical closure.
\end{proof}
\end{annexitem}

\begin{annexitem}{annex:p12.n-c-035.5.gbn-soundness}{thm:p12-eu-gbn-soundness}{Backward soundness of each exact state transformation}
\begin{proof}
If \(\rho\) literally solves solved S, then
\(\sigma_S\rho=\rho\) on the whole alphabet: each bound variable
has image \(t\rho=x\rho\), and every other variable is fixed
before \(\rho\). This calculation requires no idempotence of \(\rho\).
S itself has a literal solution \(\sigma_S\), since its domain
variables are absent from its images.

For G, backward soundness of Trivial is immediate. Decomposition's
child equalities imply the parent equality by E-congruence, including
nullary symbols. Orient uses symmetry. In Elimination a solution of
the successor literally solves the retained binding \(x\approx t\),
so \(\delta\rho=\rho\). Thus successor solution of \(P'\delta\)
modulo E implies solution of \(P'\) modulo E, and literal solution
of \(S\delta\) implies literal solution of old S. The selected
binding solves the old equation as well. In LP, successor solution
gives \(l\rho=_E u\rho\) and the fresh identity gives
\(l\rho=_E r\rho\); congruence therefore restores the old
equation \(e[u]\rho\) from its replacement \(e[r]\rho\).
Freshness and first decomposition restrict applicability but are not
needed to strengthen this backward implication.

For B/N, the first four constraint transformations preserve literal
solutions of C and S. At Basic Elimination the same
\(\delta\rho=\rho\) calculation restores C and S, while P
is unchanged. At Standard Elimination it also restores \(P\delta\)
to old P modulo E. Constrain's successor literally solves
\(e\sigma_S\); since \(\sigma_S\rho=\rho\), it solves
old e and hence its E-equation. Directed LP's successor has the
literal constraint \(l\sigma_S\rho=u\sigma_S\rho\), thus
\(l\rho=u\rho\); the rule identity and E-congruence restore
the original equation. No position inside a substitution image is
implicitly added to the B skeleton by this proof. Composite Solve
and Narrow are compositions of these sound transformations and the
literal syntactic solver, not extra unproved state operations.

Iterate these backward implications from the terminal solution
\(\sigma_S\) to the initial state. Literal constraint Clash and
Occurs have no solution by free-term roots and the finite proper-
subtree argument. Applying them to equational P would be a different,
unsound operation; identity solves \(x=_E f(x)\) in the stated
collapsing theory. Blocking only excludes the reduced witnesses used
by a completeness strategy; the successful-trace theorem makes no
claim that all E-solutions of a blocked state are absent.
\end{proof}
\end{annexitem}

\begin{annexitem}{annex:p12.n-c-035.5.cp-ground-closure}{thm:p12-eu-cp-ground-closure}{Critical closure, nested-variable peaks, and ground normal forms}
\begin{proof}
Every added critical identity is an E-consequence. At its overlap,
\(l_1|_p\mu=l_2\mu\) literally, so
\(r_1\mu=_E l_1\mu=_E(l_1[r_2]_p)\mu\).
Induction through the closure levels proves that adding these identities
does not change the E-congruence. Conversely E is included. The
oriented ground instances decrease \(\succ\), so termination follows
from its well-foundedness, even if the instance set is infinite.

Two preliminary facts justify the peak analysis. A ground term is
strictly greater than every proper subterm: otherwise totality would
give \(t\succ D[t]\) for a proper ground context D, and monotonicity
would yield the infinite descent
\(t\succ D[t]\succ D[D[t]]\succ\cdots\).
Also a nontrivial theory has no decreasing ground instance of an
identity directed from a variable. If \(x\approx r\) omits x
from r, assigning x independently of all variables in r makes every
model a singleton. If r contains x and is not x, its ground instance
properly contains that same ground image, and cannot be smaller by
the preceding context argument. Thus every actual decreasing
ground rule has a nonvariable pattern left side.

Prove by well-founded induction on a ground term w that w has a
unique normal form. Any proper reduct is smaller and already has a
unique normal form. Two first reductions from w either are disjoint
or nested. Disjoint replacements commute to a common term, so their
smaller reducts have the same normal form. For a nested overlap at a
nonvariable position of the outer left pattern, take the two fresh
pattern identities. Their actual ground match is a syntactic unifier
of the overlapping patterns, and therefore factors through a syntactic
MGU \(\mu\). The critical identity is in the next closure level.
Its ground instance is exactly the two peak branches inside their
outer context. If these branches differ, the total ground order
orients that identity from the larger branch to the smaller; if
they agree there is nothing to do. Thus their normal forms agree.

It remains to handle a nested step inside the ground image of a
variable v of the outer pattern; this is not a nonvariable critical
overlap. Write \(\gamma(v)=d\), with the inner rewrite changing
one part of d to obtain \(d'\prec d\), and let \(\gamma'\)
replace only that image. After the peak's one inner replacement,
perform the same replacement at the corresponding position in every
other occurrence of v in the outer left side. This reaches
\(l_1\gamma'\). From the outer branch \(r_1\gamma\), perform
the same replacements in all its v-images to reach \(r_1\gamma'\).
Each intermediate is below w: the left path starts with a strict
reduct of w, and the right path starts with the other strict reduct;
all further replacements decrease. The two final terms are an
instance of the outer identity, so they are identical or related by
one oriented ground identity, also entirely below w.

By induction, every term in this lower conversion has a unique
normal form, and each neighboring rewrite pair has the same one
(append a normalization of the smaller term to that rewrite).
Consequently the two peak branches have the same normal form.
This covers erasure and duplication of v and does not require the
outer identity to stay oriented in the same direction after changing
its variable image. All first reducts of w therefore normalize alike;
if w is irreducible it is its own unique normal form. This completes
the induction and proves ground confluence.

Finally, any equational proof between ground endpoints can be made
ground: its finitely many intermediate variables may all be assigned
a chosen ground term, fixing the endpoints. Each instantiated step
uses a ground identity. Such conversions have the same normal form
by confluence and termination. Conversely a common normal form gives
an equational conversion. A fresh constant supplies a ground term
if the original signature has none. No new constant occurs in a
critical identity itself, and the proof works for every finite
fresh-constant extension. This is a proof of the explicitly defined
unrestricted closure, not an identification with an unspecified
ordered-critical-pair subconstruction.
\end{proof}
\end{annexitem}

\begin{annexitem}{annex:p12.n-c-035.5.cp-elimination}{thm:p12-eu-cp-elimination}{Both source critical-pair simulations with mandatory decomposition repaired}
\begin{proof}
Give an identity its least closure depth. A finite trace uses a finite
multiset of depths. Replacing an occurrence of positive depth by
identities of smaller depth decreases the well-founded multiset order.
It is enough to simulate one use of
\(A=r_1\mu\approx B=(l_1[r_2]_p)\mu\), where the two
component identities have smaller depth and
\(\mu=\operatorname{mgu}(l_1|_p,l_2)\) is idempotent.
The chosen \(\mu\) is an actual source U output, so its recorded
finite syntactic derivation can be replayed. Choose component
variants so that variables occurring in this particular critical
identity retain its fresh names; all other component variables,
including eliminated domain variables and surviving variables absent
from the critical identity, are fresh outside the entire original
trace. Visible critical-identity variables are range variables
fixed by \(\mu\), and eliminated domains can be renamed independently.

First suppose LP replaces an occurrence u using \(A\approx B\).
Use \(r_1\approx l_1\) there, introducing \(r_1=^{?}u\)
and replacing that occurrence by \(l_1\). If this new matching
equation has nonvariable left side, immediately perform its required
first Decomposition. This is allowed: either r1 is a variable, or
its root is the root of \(r_1\mu\), which matched u in the
original admissible LP. Now use \(l_2\approx r_2\) at position
p in the newly inserted \(l_1\), producing the central constraint
\(l_2=^{?}l_1|_p\) and the replacement \(l_1[r_2]_p\).
That position is nonvariable; matching roots agree because these
patterns have an MGU. Perform its required first Decomposition
and solve the central constraint with the four nonfailure U rules.
These compute the chosen \(\mu\), or an equivalent renamed MGU;
choose the same literal solved variant by following its finite
syntactic derivation. If that derivation began with Trivial on an
identical nonvariable pair, replace that first step by Decomposition
and Trivial on its identical arguments; the resulting substitution
is still identity. The surrounding terms become
\(r_1\mu\) in the matching equation (or its already decomposed
children) and \((l_1[r_2]_p)\mu\) in the replaced context.
On all original variables the extra substitution is identity.
This is exactly the original LP successor, with harmless fresh
solved bindings adjoined. If r1 was a variable but \(r_1\mu\)
is nonvariable, now perform the matching equation's required first
Decomposition to align with the original trace's restriction.

Now suppose LP uses \(B\approx A\). A root overlap is reduced
to the first case by interchanging the two component identities:
its branches are \(r_2\mu,r_1\mu\). The interchanged central
equation may produce a different literal U output \(\nu\).
It is a syntactic MGU of the same unordered overlap problem.
The finite-permutation construction in \cref{thm:p12-ue-renaming}
gives \(\nu=\mu\pi\); both maps fix variables outside the
component patterns, so the permutation may be chosen to fix those
variables as well. Choose the fresh component range-variable names
before the simulation via \(\pi^{-1}\), retaining additional
eliminated domains as fresh names. Then its critical branch terms
are \(r_2\nu\pi^{-1}=r_2\mu\) and
\(r_1\nu\pi^{-1}=r_1\mu\). This is an explicit finite
renaming of fresh component names, not a claim that U produces
every literal idempotent MGU. For a nonroot p, first
use \(l_1\approx r_1\) at u, creating \(l_1=^{?}u\) and
inserting r1. The root of l1 is the root of B, so this step is
admissible. Immediately decompose \(l_1=^{?}u\), as required.
Because p is nonroot, its occurrence lies in the left side of one
of the resulting child equations. Paramodulate with
\(l_2\approx r_2\) at precisely that occurrence. It changes
the child of l1 to the corresponding child of \(l_1[r_2]_p\)
and introduces the central constraint
\(l_2=^{?}l_1|_p\). Decompose and solve that constraint to
compute \(\mu\). The child equations are now exactly the
first decomposition of \(B=^{?}u\), while the replacement
context is \(e[r_1\mu]\), as in the original trace.
Thus no LP is applied to an undecomposed matching equation.

The original matching equation's first Decomposition may be advanced
past operations on other equations: it introduces only its own
argument equations, and substitution commutes with decomposition
because arguments of a constructor instance are its instantiated
arguments. LP on a different occurrence is unchanged, with fresh
variants renamed away. Therefore the original continuation can be
replayed, omitting the already performed first Decomposition.
The extra solved bindings have only fresh eliminated domains;
later original eliminations update their right sides but cannot
change an original variable image merely because those bindings
are present. Original terminal bindings are retained, and the
returned map agrees on original problem variables.

Both simulations use only lower-depth identities and finite
syntactic solving. Repeating the well-founded depth-multiset
induction eliminates every positive-depth identity. The final
trace uses only E and reverse E identities. The immediate
decomposition in the second case is the explicit repair of the
printed delayed-decomposition trace, not a relaxation of LP's
earlier applicability condition.
\end{proof}
\end{annexitem}

\begin{annexitem}{annex:p12.n-c-035.5.g-completeness}{thm:p12-eu-g-completeness}{An actual ground lift and elimination back to the original G calculus}
\begin{proof}
Fix an E-solution \(\theta\) and the finite set of variables in its
relevant images. Freeze those image variables with distinct new
constants absent from E and P. Thus \(x\theta\) becomes a ground
image on each input variable; source equations and syntactic states
remain unchanged. Let R be the decreasing ground instances of the
explicit critical closure. By the preceding closure theorem it
terminates, is ground confluent in this extended signature, and
represents E. Normalize each ground image, obtaining a ground
valuation \(\alpha\) E-equal to the frozen original witness.
It remains an E-solution, and each of its variable images is R-normal.

We construct a finite successful G trace using closure identities.
At a reachable \(P;S\), keep a ground, reduced witness \(\alpha\)
that E-solves P and literally solves S. Use as induction measure
the multiset M of the ground terms occurring on both sides of P,
ordered by the well-founded multiset extension of \(\succ\),
with the number of bad nonvariable-to-variable orientations as a
second lexicographic coordinate. This is an editorial alternative
to the source's minimal-proof-length measure: the actual LP block
includes its mandatory matching Decomposition.

If some equation has literally equal instantiated sides, use the
appropriate nonfailure syntactic operation. Trivial removes two
terms from M. Decomposition replaces its two constructor terms by
proper subterms, so strictly decreases M. If the equation is
variable-to-term, proper occurrence of that variable in the term
would make its finite ground image a proper subtree of itself;
therefore the occurs proviso holds. Elimination removes that pair
from M and leaves every other instantiated term unchanged, since
\(\delta\alpha=\alpha\). It also preserves literal solution
of S. Orient leaves M unchanged and decreases the second coordinate.
All cases respect solved form. When P is empty, the construction
ends, with \(\sigma_S\alpha=\alpha\) on the input variables.

Otherwise some selected E-equal sides have different ground instances.
They have the same R-normal form, so at least one is reducible.
Choose an innermost redex. Its position is nonvariable in the current
uninstantiated term: an occurrence at or inside a substituted variable
image would be reducible inside an R-normal image, a contradiction.
Let its ground rewrite be \(l\gamma\to r\gamma\), using a
fresh closure identity with nonvariable l. The original nonvariable
subterm u and l have matching roots, and \(u\alpha=l\gamma\).
Every image of a variable occurring in l is a proper subterm of
this innermost redex and is therefore normal. Any identity variables
only in r can also be given normal images: normalize them, leaving
l's match unchanged and making r no larger. The resulting ground
instance is still decreasing. Thus extend \(\alpha\) by the
fresh reduced \(\gamma\), preserving all old values.

Apply LP at u and immediately Decompose the introduced
\(l=^{?}u\). Under the extended witness its argument equations
are literal equalities of normal proper subterms. The replaced
original ground term strictly decreases, and every newly added
argument term is smaller than its old redex, hence smaller than
the removed original term. Thus this entire admissible block
strictly decreases M. The witness remains reduced and solves all
new equations and old S. Induction now continues; every chosen
step or block decreases the stated well-founded measure.
This constructs a finite terminal trace, not just a named lifting
lemma or a assumed narrowing-success branch.

The critical-identity elimination theorem converts this finite
trace to one using E alone and preserving its answer on input
variables. Since the terminal witness literally solves S,
\(\sigma_S\alpha=\alpha\) there. Normalization changed the
original frozen witness only modulo E. Finally replace each new
parameter constant by its original variable in the finite factor
\(\alpha\). This homomorphic replacement preserves E-identities,
because no such constant occurs in E. The returned \(\sigma_S\)
contains no parameter constant: they occurred only in the witness,
not in P, closure identities, or syntactic transformations. Therefore
ungrounding yields the required original-domain comparison
\(x\theta=_E x\sigma_S\lambda\), \(x\in\mathcal X\).
This closes the dependency back to the original problem and theory.
\end{proof}
\end{annexitem}

\begin{annexitem}{annex:p12.n-c-035.5.bn-completeness}{thm:p12-eu-bn-completeness}{Reduced witnesses force actual basic positions and all induction cases}
\begin{proof}
First justify the declared grounding domain. Confluence on terms
with variables implies confluence after adding finitely many fresh
constants: replace those constants by distinct variables, use
confluence, then replace them back. Conversely, a finite variable
peak can be frozen with distinct fresh constants; a ground joining
sequence in that extension unfreezes to a joining sequence, because
the new constants occur in no rule. These homomorphisms preserve
each rule instance. Termination is supplied by the fixed reduction
ordering and is also preserved in the declared domain. Thus the
fresh-constant condition is precisely the needed convergence domain,
not a consequence of convergence on the original ground algebra.

Freeze the finite parameter set of an arbitrary solution and
normalize its input images, obtaining a reduced ground valuation
\(\alpha\) that still E-solves the initial P. The witness at a
later \(P;C;S\) must E-solve P and literally solve C and S.
Induct with the source measure
\((M,n_1,n_2,n_3)\): M is the multiset of ground terms in
\(P\alpha\), with the multiset reduction ordering;
\(n_1\) counts distinct variables in C, \(n_2\) its term
symbols, and \(n_3\) its bad orientations.

If C is nonempty, it has a literal ground solution. The four
syntactic constraint rules therefore supply an enabled nonfailure
step. Clash cannot occur between equal instantiated constructor
roots; a proper variable occurrence cannot have a finite literal
ground solution. At elimination \(\delta\alpha=\alpha\), so
C and S remain literally solved by the same reduced witness.
B leaves P unchanged; N instantiates P but \(P\delta\alpha=P\alpha\),
so M is unchanged in either system. The eliminated variable
disappears from C, decreasing n1. Trivial decreases n2 if n1
does not fall; Decomposition removes the two outer nodes if n1
does not fall; Orient lowers n3 when the earlier coordinates
are unchanged. S remains solved. These are actual enabled rules
on C, not rule choices applied to equational P.

If C is empty and an equation's ground instances are the same
normal term, Constrain removes that equation from P and places
its \(\sigma_S\)-instance in C. Because
\(\sigma_S\alpha=\alpha\), the new constraint is literally
solved by the witness. Removing its two ground terms strictly
decreases M. If P is nonempty and this case does not apply to
the selected equation, choose an innermost redex in one of its
ground sides: E-equality and confluence imply some redex unless
both sides are the same normal form.

For B the redex is at a nonvariable position of the retained
skeleton, since no reduced variable image contains a redex.
For N the same argument applies to its current instantiated
skeleton. Let \(l\gamma\to r\gamma\) be the chosen fresh
rule instance. Every \(\gamma(v)\) is a proper subterm of
the innermost redex, hence normal; the ordinary-rule variable
condition ensures there are no unconstrained right-only images.
Extend the witness by this fresh \(\gamma\). LP's added
constraint is \(l\sigma_S=^{?}u\sigma_S\); the extended
witness makes both sides literally \(l\gamma=u\alpha\),
since it literally solves S. Matching roots agree. First
Decompose the new match, as required, then use the syntactic
constraint rules. The original ground skeleton term has been
replaced by a strict reduct, so M decreases; constraints do not
contribute to this first coordinate. The witness remains reduced,
preserves old variables, and solves the successor state.

All induction cases strictly decrease the measure, and an empty
P with nonempty C is handled by the first case. The terminal
state has both empty and solved S. Its returned substitution
factors the reduced witness on input variables. Unground the
finite fresh constants as in the G proof; the source rule terms
and output contain none of those witness-only constants, and
normalization was only an E-change. This proves the required
comparison for every original solution. The same induction
proves B and N separately; it does not identify their states,
allowed positions or intermediate substitutions.

For the recorded counterexample, rule decrease in the given
positive interpretation is strict in all three cases, and all
original closed subterms normalize to a. The two-element constant-
operation model proves \(x\ne_E a\). Forward LP has only
\(g(a)\to a\) available at the root of \(g(x)\); the only
other position is the variable x, where it is prohibited.
Constrain has a constructor clash, so every successful directed
answer forces \(x=a\). Hence none covers identity. The fresh-
constant peak proves exactly why the grounding step would fail
under the weaker original-signature hypothesis. This does not
apply to G with both identity directions, which can first replace
\(g(x)\) by \(f(x,x)\) and then by a.
\end{proof}
\end{annexitem}

\begin{annexitem}{annex:p12.n-c-035.5.c-and-a}{thm:p12-eu-c-and-a}{The C-root choices and the associative word obstruction}
\begin{proof}
C-equality changes only child order at nodes of the binary symbol \(f\).
There are finitely many variants of any finite term: at most one
binary choice is needed per \(f\)-node. Every other root retains
its ordered arguments. Hence there are finitely many variant input
problems. Their successful syntactic MGUs exist by the syntactic
interface and are C-unifiers of the original problem, because
substitution preserves each input root-swap equality.

Fix a C-unifier \(\theta\). Fix a total order on finite encoded terms,
and give each term a canonical unordered-tree representative: recursively
canonicalize children, sort the two children at an \(f\)-node, and
retain argument order at other symbols. These representatives characterize
C-equality by induction on tree height. Replace each nonidentity
variable image in \(\theta\) by its canonical representative, making
a finite-support map \(\theta'\). A variable image is unchanged, and
every image of \(\theta'\) is C-equal to its image in \(\theta\).
For every original term \(s\), there is a root-swap variant \(s'\)
such that \(s'\theta'\) is the canonical representative of
\(s\theta\): at a variable this follows from the image choice;
at any root it follows by the recursive choices of the children and,
at \(f\), their one root-order choice. No inconsistency arises from
repeated variables, since all their images use the same \(\theta'\).
Thus one enumerated variant problem is syntactically solved by
\(\theta'\). Its MGU \(\mu\) factors \(\theta'\); pointwise
C-equality then gives \(x\theta=_C x\mu\lambda\) on problem
variables. This proves completeness even if images contain alien
symbols or original variables.

There are finitely many collected answers. Pass to their equivalence
classes and retain the minimal ones, taking one representative each.
Every member of a finite partial order lies above a minimal member,
so completeness is preserved. For the two-constant example any term
C-equal to \(f(a,b)\) has these two constant children in one of
the two orders. Substitution fixes constants; distinct constants are
not C-equal. The two displayed ground answers are consequently
complete and incomparable.

For A-equality, flatten a maximal \(f\)-product into a nonempty
word. Letters are variables, constants, or subterms with a different
root. Associativity only changes parentheses, so equality requires
the same outer word, with any inner theory equalities respected.
The solution condition for the displayed problem is \(aw=wa\),
where \(w\) is the nonempty flattening of \(x\theta\).
Its first letter is forced to be \(a\); cancellation reduces the
same condition to its tail. Induction on word length forces every
letter of \(w\) to be \(a\), so \(x\theta=_A a^n\), \(n\geq1\).
Conversely each such assignment solves. Different positive lengths
are not changed by substitution and are not A-equal. Therefore
\(\{\sigma_n:n\geq1\}\) is an infinite antichain complete set.
This proves that problem's type; the source's further classification
of the entire theory A depends on its separately cited general
associative procedure and is not asserted as proved by this example.
\end{proof}
\end{annexitem}

\begin{annexitem}{annex:p12.n-c-035.5.categorical-bridge}{thm:p12-eu-categorical-bridge}{Generator maps and factor arrows}
\begin{proof}
A map on generators extends uniquely by term evaluation. Every defining
identity and every congruence/substitution instance holds in the target
E-algebra, so evaluation is well defined on classes. Conversely each
homomorphism is determined by its finitely many generator images;
choosing term representatives gives a finite-support substitution,
identity outside those generators, with a finite target generator set.
Composition agrees by induction on terms. The two composite arrows
agree exactly when they agree on each \(u_i\), that is, when
\([s_i\delta]=[t_i\delta]\). An arrow factor gives finitely many
images of target generators and hence a finite substitution factor.
Conversely a substitution factor restricts to that arrow. This proves
the solution and comparison correspondences with their actual domains.

In the example every syntactic unifier sends \(x,y\) to the same
term; assigning that term to the target generator \(x\) supplies a
factor of the displayed map. Take the common-image solution \(x,y\mapsto w\)
into the free algebra on \(\{w\}\). One factor sends the unused
generator \(z\) to \(w\); another sends it to \(f(w,w)\).
They are distinct in the free term algebra but have the same composite.
Therefore the most general arrow is not a coequalizer. If its target
has only generator \(x\), that single generator value is forced,
so the factor is unique in this example. This argument does not prove
a universal target-restriction theorem for nonempty E.
\end{proof}
\end{annexitem}

\begin{annexitem}{annex:p12.n-c-035.5.coordinate-models}{thm:p12-eu-coordinate-models}{Free coordinate models and matrix composition}
\begin{proof}
For ACU flatten \(f\), remove \(e\), and reorder generator leaves.
Associativity, unit, and commutativity implement exactly these changes.
Equal multiplicity vectors therefore give theory equality. Conversely
finite-support natural-number vectors, with vector addition and zero,
model ACU; interpreting generators as unit vectors proves that theory
equality implies equal coordinates. Every finite-support vector is
represented by the indicated product, or by \(e\) if zero.
For ACUI add duplicate-leaf deletion by idempotence. Finite supports
with union form a model, proving both directions in the same way.

For AG derive inverse laws rather than presuppose them. In a commutative
monoid inverses are unique: if \(f(a,b)=e=f(a,c)\), then
\(b=f(b,e)=f(b,f(a,c))=f(f(b,a),c)=c\).
The inverse of \(e\) is \(e\), the inverse of \(i(a)\) is \(a\),
and the inverse of \(f(a,b)\) is \(f(i(a),i(b))\), by reassociation
and pair cancellation. Uniqueness yields \(i(e)=e\),
\(i(i(a))=a\), and \(i(f(a,b))=f(i(a),i(b))\).
Push inverses to generator leaves, cancel opposite occurrences, and
collect the remaining signed counts. Equal coordinates thus imply
equality; the finite-support integer vector group is a model, proving
the converse. Its negative coordinates are represented by inverse
leaves, and zero by \(e\).

For a substitution, the coordinate at \(q\) of a substituted term
is the sum over input generators \(x\) of their original coordinates
times the coordinates at \(q\) of \(x\sigma\). This is induction
over the constructors, using the appropriate semiring addition and
multiplication; the AG inverse case uses negation. Applying this
calculation twice gives the source-order matrix product. Added free
constants are generator leaves fixed by every substitution, so this
proof also applies to problems with constants.
\end{proof}
\end{annexitem}

\begin{annexitem}{annex:p12.n-c-035.5.homogeneous-completeness}{thm:p12-eu-homogeneous-completeness}{All solution columns and finite generation}
\begin{proof}
Each fresh-variable column of \(\sigma\) is a homogeneous solution,
so the coordinate criterion proves soundness. For any unifier
\(\theta\), let \(Q\) be the finite set of generators occurring
in its images. Each coordinate column
\(w_q=(\operatorname{coef}_q(x_i\theta))_i\) is in \(U\), and
write \(w_q=\sum_jv_js_{j,q}\). Assign to \(z_j\) the finite
term having coefficients \(s_{j,q}\) at all \(q\in Q\).
This gives a finite-support factor \(\lambda\), allowing zero
terms and, for AG, inverse terms. Its substituted coordinate at \(q\)
is \(\sum_j(v_j)_is_{j,q}=\operatorname{coef}_q(x_i\theta)\).
All other coordinates vanish. Hence
\(x_i\theta=_E x_i\sigma\lambda\) on the input variables.
Fresh basis variables are outside the input; \(Q\) may contain those
variables without affecting this simultaneous calculation. Comparison
is on input variables, not on the unused alphabet.

Here is the required finite-generation argument for ACU. Every infinite
sequence of natural numbers has an infinite nondecreasing subsequence:
either a value repeats infinitely, or one can choose successively larger
values because each bounded range has only finitely many occurrences.
Apply this successively to the \(n\) coordinates of an infinite
sequence in \(\mathbb N^n\). The resulting subsequence contains a
componentwise comparable pair in their original index order. Therefore
there is no infinite antichain, and the minimal elements of any subset
are finite. For any nonzero homogeneous solution \(w\), its finite
box contains a minimal nonzero solution \(v\leq w\). Subtraction
gives a nonnegative homogeneous solution \(w-v\). Induction on the
sum of coordinates expresses every \(w\) as a sum of minimal
nonzero solutions; zero is the empty sum. Those minimal solutions form
an antichain and hence are finite. This proves existence and generation,
not a bounded stopping rule for a search for them. ACUI has finitely
many solution vectors in \(BS^n\), and each is generated by itself
with scalar one.

For clarity, integer homogeneous kernels also have finite generators.
Every subgroup \(H\) of \(\mathbb Z^n\) is finitely generated,
by induction on \(n\). For \(n=0\) it is zero. Project onto the
last coordinate. If the image is zero, apply the induction to the
first \(n-1\) coordinates. Otherwise its positive values have a
least \(d\); integer division shows the image is \(d\mathbb Z\).
Choose \(h\in H\) with last coordinate \(d\), and generate the
kernel of that projection by induction. For every \(w\in H\)
with last coordinate \(kd\), the difference \(w-kh\) lies in
that kernel. Adjoining \(h\) therefore generates \(H\). The
integer solution set of the homogeneous matrix equality is a subgroup,
so this applies. It establishes finiteness and the factorization used
here; effective integer normal-form algorithms remain the separate
computational methods cited by the source.
\end{proof}
\end{annexitem}

\begin{annexitem}{annex:p12.n-c-035.5.constant-completeness}{thm:p12-eu-constant-completeness}{Fixed constants and the particular-column construction}
\begin{proof}
Each chosen constant column satisfies its inhomogeneous equality and
each fresh-variable column its homogeneous equality. The coordinate
criterion proves soundness. For an arbitrary unifier \(\theta\),
each constant column is \(p_c+v_c\), with \(p_c\in P_c\) and
\(v_c\in U\). Select that tuple of particular solutions, producing
one output substitution. Generate all \(v_c\)'s and all variable
columns of \(\theta\) from the homogeneous generators. As in the
previous annex, assign those scalar coefficients at constants and image
variables to the fresh basis variables. Composition adds each fixed
\(p_c\) to its generated \(v_c\), and reproduces every coordinate
of \(\theta\). Constants themselves never receive new images.
This is the required comparison-domain factorization. There are finitely
many choices of particular columns, so the set is finite complete.
An empty inhomogeneous solution set rules out any unifier, which would
otherwise supply such a column.

In \(\mathbb N^n\), minimal inhomogeneous solutions form a finite
antichain by the preceding Dickson argument. For any solution \(w\),
choose a minimal solution \(p\leq w\) in its finite box. Subtracting
their inhomogeneous integer equalities gives the homogeneous equality
for the nonnegative vector \(w-p\). This proves the particular-set
coverage property. In \(BS^n\) simply take all inhomogeneous solutions;
each is itself plus homogeneous zero, and the set is finite.
Over \(\mathbb Z\), one solution \(p\) gives all solutions as
\(p+v\), \(v\in U\), by subtraction and addition. Thus one
particular tuple gives a single complete answer for a solvable AG
problem with constants.

For the body's ACU constant example, every solution of
\(2x+y=3z+1\) with \(y>0\) is \((0,1,0)\) plus a
nonnegative homogeneous vector. If \(y=0\), then
\(2x=3z+1\) forces \(x=2+3k,z=1+2k\), \(k\geq0\),
so it is \((2,0,1)\) plus a homogeneous vector. These two
particular vectors are componentwise incomparable and minimal.
The body's constant-occurrence arguments prove incomparability of the
two resulting substitutions; together with completeness they prove
minimality. ACUI can also require several answers: for free \(a\),
\(f(x,y)=^{?}_{ACUI}a\) has exactly the three ground answers
\((x,y)=(a,e),(e,a),(a,a)\). Each image support is a subset of
\(\{a\}\), and the two supports must have union \(\{a\}\).
These three possibilities are exhaustive and, being distinct ground
assignments, incomparable. This proves the stated nonsingleton
qualification rather than deriving it from merely redundant outputs.
\end{proof}
\end{annexitem}

\begin{annexitem}{annex:p12.n-c-035.5.semiring-criterion}{thm:p12-eu-semiring-criterion}{Finite matrix factors with fixed-generator blocks}
\begin{proof}
In a semi-additive category a map between finite biproducts is
determined by its one-generator components. Composition sums through
intermediate components, giving matrix multiplication; every finite
matrix of endomorphism entries represents an arrow. These are precisely
the matrix hypotheses in the statement, not extra properties of an
arbitrary free-function extension.

If \(U\) has generators \(v_1,\ldots,v_k\), let \(H\) have
those columns. It satisfies \(M_1H=M_2H\). Every other solution
matrix \(D\) has all columns in \(U\); putting their generating
scalars into a matrix \(L\) gives \(D=HL\), proving most
generality. For \(k=0\), use the zero object as target; every
solution column is zero and factors through it. Conversely a most
general arrow has a finite target generator set, and hence finitely
many columns in its matrix \(H\). Every \(v\in U\) is a
single-column solution arrow into the one-generator free algebra and
factors through \(H\). Therefore \(v=H\ell\), proving that
the columns generate \(U\). The zero arrow ensures elementary
solvability under the standing semi-additive hypothesis.

Order variable generators before constant generators. Direct block
multiplication gives
\[
 \begin{pmatrix}A&B\\0&I\end{pmatrix}
 \begin{pmatrix}H&P\\0&I\end{pmatrix}
 =\begin{pmatrix}AH&AP+B\\0&I\end{pmatrix},
\]
proving the two conditions in the body. Suppose finitely many
particular translates cover every constant-column solution set, and
let \(H\) generate homogeneous columns. Each finite tuple \(P\)
gives a solution output. An arbitrary solution factors through an
output: its variable block is \(HL\), and its constant block is
\(P+HK\), exactly the product blocks of a constant-respecting
factor \(\left(\begin{smallmatrix}L&K\\0&I\end{smallmatrix}\right)\).
This proves a finite complete set.

Conversely, for given matrices and vectors form a parallel pair with
one fixed generator \(c\), whose constant columns are \(u_1,u_2\).
A finite complete set with constants has homogeneous blocks
\(H_1,\ldots,H_k\) and constant columns \(p_1,\ldots,p_k\).
Any inhomogeneous solution \(w\) is a ground answer matrix,
with that constant column and no fresh-variable columns. Factorization
through some output gives \(w=p_j+H_j\ell\); all columns of
\(H_j\) are homogeneous, hence \(H_j\ell\in U\).
Conversely any \(p_j+v\), \(v\in U\), solves by distributivity.
Thus the inhomogeneous set is exactly the finite union of those
translates. The empty solution set is represented by the empty union.
This establishes the equivalence under the stated matrix interface;
it supplies no normalization or completeness proof for combination.
\end{proof}
\end{annexitem}
